
\begin{tikzpicture}[scale=1.2]

% 점 정의
\tkzDefPoint(0,0){O}
\tkzDefPoint(-128:3){B}
\tkzDefPoint(-52:3){C}
\tkzDefPoint(68:3){A}

\tkzDrawPoint(O)    \tkzLabelPoint[below](O){$O$}
\tkzDrawPoint(A)    \tkzLabelPoint[above](A){$A$}
\tkzDrawPoint(B)    \tkzLabelPoint[below left](B){$B$}
\tkzDrawPoint(C)    \tkzLabelPoint[below right](C){$C$}
\tkzDrawCircle[dashed](O,A)

\tkzDrawSegments(A,B B,C C,A)

\tkzDefCircle[circum](B,C,O) \tkzGetPoint{X}
\tkzDrawCircle(X,O)

\tkzInterLC(A,B)(X,O) \tkzGetPoints{D}{t}
\tkzDrawPoint(D)    \tkzLabelPoint[left](D){$D$}

\tkzInterLC(A,C)(X,O) \tkzGetPoints{t}{E}
\tkzDrawPoint(E)    \tkzLabelPoint[right](E){$E$}

\tkzInterLL(D,O)(A,C) \tkzGetPoint{P}
\tkzDrawPoint(P)    \tkzLabelPoint[right](P){$P$}

\tkzInterLL(E,O)(A,B) \tkzGetPoint{Q}
\tkzDrawPoint(Q)    \tkzLabelPoint[above left ](Q){$Q$}
\tkzDrawSegment(E,Q)
\tkzDrawSegment[dashed](O,B)

% 반원 (AB가 지름)
% \tkzDrawSemiCircle[thick](O,B)

% % 지름 AB
% \tkzDrawSegment[thick](A,B)

% % CD를 지름으로 하는 원
% \tkzDrawCircle[thick, color=blue!60](M,C)

% % 선분 CD
% \tkzDrawSegment[thick, color=red!70](C,D)

% % 접점 E에서 수직선 (접선의 법선)
% \tkzDrawSegment[dashed, color=gray](E,M)

% % 선분 OE 표시
% \tkzDrawSegment[thick, color=green!50!black](O,E)

% % 점 표시
% \tkzDrawPoints[size=4](A,B,O,E,M,C,D)

% % 라벨
% \tkzLabelPoint[below left](A){$A$}
% \tkzLabelPoint[below right](B){$B$}
% \tkzLabelPoint[below](O){$O$}
% \tkzLabelPoint[below](E){$E$}
% \tkzLabelPoint[right](M){$M$}
% \tkzLabelPoint[above left](C){$C$}
% \tkzLabelPoint[above right](D){$D$}

% % 길이 표시
% \tkzLabelSegment[below, pos=0.5](O,E){1}
% \tkzLabelSegment[above, pos=0.5, color=red!70](C,D){12}

% % 직각 표시 (접점에서)
% \tkzMarkRightAngle[size=0.4](M,E,B)

\end{tikzpicture}